% titel: Blatt 3, A8: Wärmekapazität eines Kalorimeters
% Teil II: der durchgerechnete Weg.

\subsection*{Gegeben und gesucht}
$Q = \qty{25.00}{kJ}$ zugefügt; Temperatur steigt von
$T_{\text{anf}} = \qty{19.84}{\celsius}$ auf
$T_{\text{end}} = \qty{21.12}{\celsius}$. Gesucht: Wärmekapazität
$C$ des Kalorimeters.

\subsection*{Ansatz}
$C = Q/\Delta T$ mit $\Delta T = T_{\text{end}} - T_{\text{anf}}$. Die
gegebenen Werte $c_p(\mathrm{H_2O})$ und $\rho(\mathrm{H_2O})$ werden nicht
gebraucht, weil keine Wassermasse in der Aufgabe steckt.

\subsection*{Rechnung}
\begin{align*}
  \Delta T &= 21.12 - 19.84 = \qty{1.28}{K}, \\
  C &= \frac{25000\ \mathrm{J}}{1.28\ \mathrm{K}}
  \approx \qty{19531}{J/K}
  \approx \qty{1.95e4}{J/K}
  = \qty{19.5}{kJ/K}.
\end{align*}

\begin{ergebnis}
$C \approx \qty{1.95e4}{J/K}$ bzw.\ $C \approx \qty{19.5}{kJ/K}$.
\end{ergebnis}

\begin{kontrolle}
Grössenordnung: eine Wärmekapazität von $\approx 20\ \mathrm{kJ/K}$ steckt
typischerweise einige Liter Wasser (über $m c_p \approx
20\ \mathrm{kJ/K}$); mit $Q = C\,\Delta T$ zurück gerechnet liefert
$19531\times 1.28 = \qty{25000}{J}$, also die vorgegebene Energie. Die
Wasser-Konstanten ($c_p$, $\rho$) ändern das Resultat nicht, da keine
Wassermasse gegeben ist.
\end{kontrolle}

\begin{merksatz}
Aus zugeführter Wärme und beobachteter Temperaturerhöhung folgt die
Wärmekapazität direkt über $C = Q/\Delta T$; die Differenz in °C oder K ist
identisch.
\end{merksatz}
